% !TEX root = ../main.tex

\begin{titlepage}
\thispagestyle{empty}
\begin{tikzpicture}[remember picture,overlay]
  \fill[techblue!5] (current page.south west) rectangle (current page.north east);
  \fill[techblue] (current page.south west) rectangle ([xshift=1.05cm]current page.north west);
  \fill[techorange] ([xshift=1.05cm]current page.south west) rectangle ([xshift=1.20cm]current page.north west);
  \begin{scope}[shift={([xshift=-3.1cm,yshift=-3.0cm]current page.north east)}]
    \foreach \rx/\ry/\ang in {0.80/0.24/18,1.20/0.36/-22,1.65/0.48/39,2.10/0.62/-8,2.65/0.78/26} {
      \draw[techblue!24,line width=0.45pt,rotate=\ang] (0,0) ellipse [x radius=\rx cm, y radius=\ry cm];
    }
    \fill[techblue!55] (0,0) circle (1.2pt);
    \foreach \x/\y in {-0.7/0.5,0.9/0.4,-0.4/-0.7,0.55/-0.45,1.25/-0.15} {
      \fill[techorange!70] (\x,\y) circle (0.8pt);
    }
  \end{scope}
  \draw[techblue!18,line width=0.6pt] ([xshift=1.65cm,yshift=-2.0cm]current page.north west) -- ([xshift=-1.65cm,yshift=-2.0cm]current page.north east);
  \draw[techblue!18,line width=0.6pt] ([xshift=1.65cm,yshift=2.0cm]current page.south west) -- ([xshift=-1.65cm,yshift=2.0cm]current page.south east);
\end{tikzpicture}

\vspace*{1.4cm}
\hspace*{1.45cm}
\begin{minipage}{0.78\textwidth}
\raggedright
{\sffamily\color{techgray}\Large 技术综述}\\[1.1cm]

{\fontsize{30pt}{38pt}\selectfont\bfseries\color{techblue}
可信 3D Gaussian Splatting\\[0.25em]
方法综述}\\[0.8cm]

{\Large\color{techgray}
从 3DGS 锚点到可验证技术路线}\\[1.2cm]

{\color{techblue}\rule{0.56\textwidth}{0.8pt}}\\[1.0cm]

{\large
934 篇方法 \quad 21 章技术路线 \quad 逐方法精写}\\[0.45cm]
{\normalsize\color{techgray}
企业背景、代码链接、GitHub 关注度与引用信息随文说明}\\[2.2cm]

{\large Codex 逐方法校订版}\\[0.35cm]
{\normalsize 2026 年 7 月版}
\end{minipage}

\vfill
\hspace*{1.45cm}
\begin{minipage}{0.78\textwidth}
\raggedright
{\small\color{techgray}
本书以“问题背景、技术路线、锚点方法、增量方法”为组织方式，关注每个方法相对原始 3DGS 或局部锚点具体改了什么、如何训练和渲染、哪些外部信息已经可以核验。}
\end{minipage}
\vspace*{1.0cm}
\end{titlepage}

\clearpage
\thispagestyle{empty}
\vspace*{0.16\textheight}

{\Large\bfseries 版本说明}\\[1.2cm]

\noindent\begin{tabularx}{\textwidth}{@{}p{3.2cm}X@{}}
\textbf{书名} & 可信 3D Gaussian Splatting 方法综述 \\
\textbf{副标题} & 从 3DGS 锚点到可验证技术路线 \\
\textbf{版本} & Codex 逐方法校订版，2026 年 7 月版 \\
\textbf{正文范围} & 收录与 3DGS 高相关、且至少具备一项可核验外部信息的方法。低相关工作保留在数据库，不进入本书正文。 \\
\textbf{写作原则} & 不做排行榜，不把外部信息写成质量分。正文以技术解释为主，外部信息只帮助读者判断方法是否值得进一步阅读、复现或跟踪。 \\
\textbf{持续更新} & 代码状态、GitHub stars 与引用数会变化；本书保留待跟进清单，后续更新时按外部信息更新流程和 Codex 逐方法审校继续增量修订。 \\
\end{tabularx}

\vfill
{\small\color{techgray}
本书为研究与工程参考材料。文中对论文、代码和引用的描述来自本地论文数据库、论文源码、已抓取代码材料和审校时可核验的信息；新论文的公开代码与引用状态需要持续复核。}

\clearpage
